
\begin{frame}[fragile]
	\frametitle{État de l'art}
	\centering
	\tiny
	\begin{tblr}{
		colspec = {|m{2.5cm}|X[c,m]|X[c,m]|X[c,m]|X[c,m]|X[c,m]|X[c,m]|},
		% cell{0}{0-} = {cmd=\SetCell{wd=5cm}},
		cell{2}{2} = {c=4}{c}, % multicolumn
		% rowsep = 6pt,
		hlines={black,0.5pt},
		vlines={black,0.5pt},
		% vlines = {red5, 1pt}, % vlines can not pass through multicolumn cells
		}
		% \bottomrule%
		\diagbox[font=\tiny,width=\dimexpr\linewidth+2\tabcolsep,height=\dimexpr+2\line+\abovesep+\belowsep+2pt]{Approche}{Système}
		                                                                                & Kripac                                                             & Capoyleas et al. & Wu et al. & Bidarra et al. & Cardot & Cascaval
		et al.                                                                                                                                                                                                                   \\ %
		\toprule[black,0.5pt] %
		% \bottomrule %
		Types d'éditions~:                                                              & paramètres géométriques uniquement                                 &                  &           &                &
		paramètres géométriques\par + \par ajout, suppression, déplacement d'opérations & paramètres géométriques\par + \par ajout, suppression d'opérations                                                                     \\ %

		Historique d'une cellule~:
		\par \textbullet~construction (complet / partiel)
		\par \textbullet~calcul (local / global)                                        &
		\par ~
		\par partiel
		\par ~
		\par global                                                                     &
		\par ~
		\par partiel
		\par ~
		\par global                                                                     &
		\par ~
		\par partiel
		\par ~
		\par local                                                                      &
		\par ~
		\par partiel
		\par ~
		\par global                                                                     &
		\par ~
		\par partiel
		\par ~
		\par local                                                                      &
		\par ~
		\par partiel
		\par ~
		\par local                                                                                                                                                                                                               \\ %
		Nomination~:
		\par \textbullet~base
		\par \textbullet~homogène                                                       &
		\par ~
		\par face
		\par \xmark                                                                     &
		\par ~
		\par face
		\par \xmark                                                                     &
		\par ~
		\par face
		\par \xmark                                                                     &
		\par ~
		\par face
		\par \xmark                                                                     &
		\par ~
		\par brin
		\par \cmark                                                                     &
		\par ~
		\par contraintes
		\par \cmark                                                                                                                                                                                                              \\ %

		Dimension du modèle                                                             & 3D                                                                 & 3D               & 3D
		                                                                                & 3D                                                                 & \(n\)-D          & 2.5D                                           \\ %
		Opérations complexes                                                            &
		\cmark                                                                          & \cmark                                                             & \cmark           & \cmark    & \xmark         & \cmark            \\ %
	\end{tblr}
	\begin{tikzpicture}[overlay, remember picture]
		\node at (current page.north)[below = 1.2cm]
		{%
			\only<2>{%
				\begin{minipage}[t][\textheight][s]{.4\textwidth}
					\begin{myblock}
						\includegraphics[width=\textwidth]{types-editions-operations-cardot}
					\end{myblock}
				\end{minipage}
			}
		};
		\node at (current page.north)[below = 1.2cm]
		{%
			\only<4>{%
				\begin{minipage}[t][\textheight][s]{.4\textwidth}
					\begin{myblock}[]
						\includegraphics[width=\textwidth]{kripac-faceid-graph}
						\includegraphics[width=\textwidth]{wu-exemple-propagation-fusion}
					\end{myblock}
				\end{minipage}
			}%
		};
		\node at (current page.north)[below = 1.2cm]
		{%
			\only<6>{%
				\begin{minipage}[t][\textheight][s]{.4\textwidth}
					\begin{myblock}
						\centering
						\includegraphics[width=.9\textwidth]{bidarra-exemple-nomination-arete-1}
						\includegraphics[width=.9\textwidth]{bidarra-exemple-nomination-arete-2}
					\end{myblock}
				\end{minipage}
			}
		};
		\node at (current page.north)[below  = 1.2cm]
		{%
			\only<8>{%
				\begin{minipage}[t][\textheight][s]{.6\textwidth}
					\begin{myblock}
						\includegraphics[width=\textwidth]{cascaval-2-5D-examples}
					\end{myblock}
				\end{minipage}
			}
		};
		\node at (current page.north)[below = 1.2cm]
		{%
			\only<10>{%
				\begin{minipage}[t][\textheight][s]{.6\textwidth}
					\begin{myblock}
						\includegraphics<10>[width=\textwidth]{wu-diff-operation-examples}
					\end{myblock}
				\end{minipage}
			}
		};
	\end{tikzpicture}
\end{frame}


%

\begin{frame}
	\frametitle{Limitations des approches existantes}
	\begin{itemize}[<+->]
		\item édition limitée aux paramètres géométriques
		\item manque d'homogénéité dans la nomination
		\item suivi partiel des cellules
		\item conception généralement limitée à la 3D
	\end{itemize}
\end{frame}

\begin{frame}
	\frametitle{Objectifs}
	\begin{columns}[onlytextwidth]
		\begin{column}{.6\textwidth}
			\begin{itemize}[<+->]
				\item Édition par réorganisation des opérations (ajout, suppression, déplacement)
				\item Schéma de nomination homogène pour tout type de cellule
				\item Suivi complet de l'évolution des cellules
				\item Modélisation d'objet générique en dimension
			\end{itemize}
		\end{column}
		\begin{column}<+->{.4\textwidth}
			\vspace{5.1em}%
			$\rdelim\}{4.5}[-4.7em]{-5em}[]$%
			\tabto{1em}formalisation des\par
			\tabto{1em}objets\par
			\tabto{1em}opérations\par
			\tabto{1em}changements topologiques%
		\end{column}
	\end{columns}
\end{frame}


\begin{frame}
	\frametitle{Objectifs}
		\begin{itemize}
			\item Conception d'un schéma de nomination homogène
			      \rdelim\}{5.5}{*}[Objectifs]
			\item Suivi complet des évolutions de cellules
			\item Édition étendue de constructions (ajout,\par suppression
			      et déplacement d'opérations)
			\item Rejeu de constructions en toutes dimensions
		\end{itemize}

	\vspace{1em}

		\begin{itemize}
			\item Cartes généralisées : formalisation des objets
			      \rdelim\}{2.2}{*}[Formalismes]
			\item Règles Jerboa : formalisation des opérations
		\end{itemize}

	\vspace{1em}

		\begin{itemize}
			\item Méthodes de détection des changements topologiques
			      \rdelim\}{3.5}{*}[Contributions]
			\item Mécanisme de rejeu basé sur des règles Jerboa
			\item Extension pour le support d'opérations complexes
		\end{itemize}
\end{frame}




%%% Local Variables:
%%% mode: LaTeX
%%% TeX-master: "../../main"
%%% End:
